\input format.tex

\usepackage{graphicx}
\graphicspath{{cores/}}

\usepackage{environ}
\usepackage{colortbl,array,booktabs}
\usepackage{tabularx}

\colorlet{TablaBordeSuperior}{topcolor}
\colorlet{TablaBordeInferior}{topcolor}
\colorlet{TablaCentroSuperior}{blue!1}
\colorlet{TablaCentroInferior}{blue!20}
\colorlet{FuenteCabeceraTabla}{white}

\newcolumntype{M}[1]{>{\centering\arraybackslash}m{#1}}

\tcbset{rtab/.style={
freelance,
frame code={
 \path[top color=topcolor,bottom color=topcolor]
   ([yshift=-#1*(\baselineskip+2pt)]interior.north west) --
   ([yshift=-#1*(\baselineskip+2pt)]interior.north east) {[sharp corners]--
    ([yshift=3pt]interior.north east) --
    ([yshift=3pt]interior.north west)} -- cycle;

  },
interior code={},
 }
}

\newcommand\fuentecabecera[1]{\textcolor{black}{\textbf{#1}}}

\NewEnviron{RCtable}[1]
  {%
    \begin{tcolorbox}[#1]\BODY\end{tcolorbox}%
  }



\begin{document}

\vspace*{8mm}
%% 各章节
\setlength{\arrayrulewidth}{1pt}
\fontsize{9.3pt}{11pt}\selectfont
\color{gray2}

\begin{center}
{\noindent\bf\sanhao 检测结果}
\end{center}

\vspace*{4mm}
\noindent
{\noindent\bf\xiaosihao 1、肠道菌群概况}
\vspace*{1.5mm}

%\begin{RCtable}{
%  rtab=1,
%  tabularx*={\arrayrulecolor{gray2}}%
%    {M{0.3\linewidth} M{0.31\linewidth}M{0.34\linewidth}}
%}
%\fuentecabecera{测试项目} & \fuentecabecera{健康状态提示} & \fuentecabecera{检测结果评价} \\[3pt]
%\bf{检测项目} & \bf{健康状态提示} & \bf{检测结果评价} \\[3pt]
%\end{RCtable}

\begin{tctabularx}{tabularx={m{4.8cm}<{\centering}m{5.2cm}<{\centering}m{4.55cm}<{\centering}}}
&&
\\[-6pt]
\cellcolor{topcolor} \raisebox{2.614pt}{\color{black00}\bf 检测项目} &
\cellcolor{topcolor} \raisebox{2.614pt}{\color{black00}\bf 健康状态提示} &
\cellcolor{topcolor} \raisebox{2.614pt}{\color{black00}\bf 检测结果评价}
\end{tctabularx}

\vspace*{-4.25mm}
\fontsize{8.8pt}{11pt}\selectfont
\arrayrulecolor{gray2}
\begin{longtable}{m{4.8cm}m{5.2cm}<{\centering}m{0cm}@{}m{4.61cm}<{\centering}}
\hline
\parbox[c]{\hsize}{\vskip7pt {\lantxh 肠道菌群多样性\\用于衡量肠道菌群物种丰富度及稳定性的指标} \vskip7pt} & \parbox[c]{\hsize}{\vskip7pt\centerline{\raisebox{-1.5ex}{\includegraphics[width=5cm,height=0.55cm]{bar.pdf}}}\vskip7pt}  &
\hspace*{-4.83cm}\raisebox{-0.45ex}{\includegraphics{smile.pdf}}
 & \begin{minipage}{4.60cm}\begin{center}{{\lantxh 正常{\\ \bahao 肠道菌群组成较为丰富，菌群失调风险较低。}} }\end{center} \end{minipage} \\
\hline
\parbox[c]{\hsize}{\vskip7pt {\lantxh 有益菌\\产生有益物质，抑制病原菌的生长，帮助预防疾病} \vskip7pt} & \parbox[c]{\hsize}{\vskip7pt\centerline{\raisebox{-1.5ex}{\includegraphics[width=5cm,height=0.55cm]{bar.pdf}}}\vskip7pt}  &
\hspace*{-4.83cm}\raisebox{-0.45ex}{\includegraphics{smile.pdf}}
 & \begin{minipage}{4.60cm}\begin{center}{{\lantxh 偏高{\\ \bahao 产生较多有益物质，帮助抑制致病菌及肠道炎症，维持肠道健康}} }\end{center} \end{minipage} \\
\hline
\parbox[c]{\hsize}{\vskip7pt {\lantxh 有害菌\\产生多种毒素，危害肠道健康} \vskip7pt} & \parbox[c]{\hsize}{\vskip7pt\centerline{\raisebox{-1.5ex}{\includegraphics[width=5cm,height=0.55cm]{bar.pdf}}}\vskip7pt}  &
\hspace*{-4.83cm}\raisebox{-0.45ex}{\includegraphics{smile.pdf}}
 & \begin{minipage}{4.60cm}\begin{center}{{\lantxh 正常{}} }\end{center} \end{minipage} \\
\hline
\end{longtable}

\vspace*{2mm}
\noindent
{\noindent\bf\xiaosihao 2、常见肠道菌含量}
\vspace*{1.5mm}

\begin{tctabularx}{tabularx={m{4.8cm}<{\centering}m{5.2cm}<{\centering}m{4.55cm}<{\centering}}}
&&
\\[-6pt]
\cellcolor{topcolor} \raisebox{2.614pt}{\color{black00}\bf 检测项目} &
\cellcolor{topcolor} \raisebox{2.614pt}{\color{black00}\bf 健康状态提示} &
\cellcolor{topcolor} \raisebox{2.614pt}{\color{black00}\bf 检测结果评价}
\end{tctabularx}

\vspace*{-4.25mm}
\fontsize{8.8pt}{11pt}\selectfont
\arrayrulecolor{gray2}
\begin{longtable}{m{4.8cm}m{5.2cm}<{\centering}m{0cm}@{}m{4.61cm}<{\centering}}
\hline
\parbox[c]{\hsize}{\vskip7pt {\lantxh 链球菌属\\多数为共生菌，少数在免疫力弱时可能引起化脓性炎症、肺炎等} \vskip7pt} & \parbox[c]{\hsize}{\vskip7pt\centerline{\raisebox{-1.5ex}{\includegraphics[width=5cm,height=0.55cm]{bar.pdf}}}\vskip7pt}  &
\hspace*{-1.51cm}\raisebox{-0.45ex}{\includegraphics{cry.pdf}}
 & \begin{minipage}{4.60cm}\begin{center}{{\color{red}\lantxh 高{\\ \bahao 可能增加化脓性炎症、肺炎等的风险}} }\end{center} \end{minipage} \\
\hline
\parbox[c]{\hsize}{\vskip7pt {\lantxh 弯曲杆菌属\\多数菌种为致病菌，可引起弯曲菌病，表现为严重腹泻或痢疾综合征} \vskip7pt} & \parbox[c]{\hsize}{\vskip7pt\centerline{\raisebox{-1.5ex}{\includegraphics[width=5cm,height=0.55cm]{bar.pdf}}}\vskip7pt}  &
\hspace*{-3.17cm}\raisebox{-0.45ex}{\includegraphics{cry.pdf}}
 & \begin{minipage}{4.60cm}\begin{center}{{\lantxh 偏高{\\ \bahao 可能增加腹泻或痢疾风险}} }\end{center} \end{minipage} \\
\hline
\parbox[c]{\hsize}{\vskip7pt {\lantxh 栖粪杆菌属\\发酵纤维素产生丁酸等有益物质，抑制肠道炎症，促进肠道健康} \vskip7pt} & \parbox[c]{\hsize}{\vskip7pt\centerline{\raisebox{-1.5ex}{\includegraphics[width=5cm,height=0.55cm]{bar.pdf}}}\vskip7pt}  &
\hspace*{-4.83cm}\raisebox{-0.45ex}{\includegraphics{smile.pdf}}
 & \begin{minipage}{4.60cm}\begin{center}{{\lantxh 高{\\ \bahao 有利于产生有益物质及抑制肠道炎症}} }\end{center} \end{minipage} \\
\hline
\end{longtable}

\vspace*{2mm}
\noindent
{\noindent\bf\xiaosihao 3、肠道致病菌}
\vspace*{1.5mm}

\begin{tctabularx}{tabularx={m{4.8cm}<{\centering}m{5.2cm}<{\centering}m{4.55cm}<{\centering}}}
&&
\\[-6pt]
\cellcolor{topcolor} \raisebox{2.614pt}{\color{black00}\bf 检测项目} &
\cellcolor{topcolor} \raisebox{2.614pt}{\color{black00}\bf 健康状态提示} &
\cellcolor{topcolor} \raisebox{2.614pt}{\color{black00}\bf 检测结果评价}
\end{tctabularx}

\vspace*{-4.25mm}
\fontsize{8.8pt}{11pt}\selectfont
\arrayrulecolor{gray2}
\begin{longtable}{m{4.8cm}m{5.2cm}<{\centering}m{0cm}@{}m{4.61cm}<{\centering}}
\hline
\parbox[c]{\hsize}{\vskip7pt {\lantxh 龋齿罗斯氏菌\\可能引起口腔疾病和其他部位的感染} \vskip7pt} & \parbox[c]{\hsize}{\vskip7pt\centerline{\raisebox{-1.5ex}{\includegraphics[width=5cm,height=0.55cm]{bar.pdf}}}\vskip7pt}  &
\hspace*{-3.17cm}\raisebox{-0.45ex}{\includegraphics{cry.pdf}}
 & \begin{minipage}{4.60cm}\begin{center}{{\lantxh 偏高{\\ \bahao 可能增加口腔疾病、感染的风险}} }\end{center} \end{minipage} \\
\hline
\parbox[c]{\hsize}{\vskip7pt {\lantxh 非解乳糖链球菌\\可导致新生儿败血症、心内膜炎} \vskip7pt} & \parbox[c]{\hsize}{\vskip7pt\centerline{\raisebox{-1.5ex}{\includegraphics[width=5cm,height=0.55cm]{bar.pdf}}}\vskip7pt}  &
\hspace*{-3.17cm}\raisebox{-0.45ex}{\includegraphics{cry.pdf}}
 & \begin{minipage}{4.60cm}\begin{center}{{\lantxh 偏高{\\ \bahao 可能增加新生儿败血症、心内膜炎的风险}} }\end{center} \end{minipage} \\
\hline
\parbox[c]{\hsize}{\vskip7pt {\lantxh 咽峡炎链球菌\\属于人体正常菌群一部分，但也可能造成肝脓肿、脑脓肿、菌血症等} \vskip7pt} & \parbox[c]{\hsize}{\vskip7pt\centerline{\raisebox{-1.5ex}{\includegraphics[width=5cm,height=0.55cm]{bar.pdf}}}\vskip7pt}  &
\hspace*{-3.17cm}\raisebox{-0.45ex}{\includegraphics{cry.pdf}}
 & \begin{minipage}{4.60cm}\begin{center}{{\lantxh 偏高{\\ \bahao 可能增加肝脓肿、脑脓肿、菌血症等的风险}} }\end{center} \end{minipage} \\
\hline
\parbox[c]{\hsize}{\vskip7pt {\lantxh 脆弱拟杆菌\\可能导致菌血症、腹内感染、腹膜炎} \vskip7pt} & \parbox[c]{\hsize}{\vskip7pt\centerline{\raisebox{-1.5ex}{\includegraphics[width=5cm,height=0.55cm]{bar.pdf}}}\vskip7pt}  &
\hspace*{-4.83cm}\raisebox{-0.45ex}{\includegraphics{smile.pdf}}
 & \begin{minipage}{4.60cm}\begin{center}{{\lantxh 含量未超标{}} }\end{center} \end{minipage} \\
\hline
\parbox[c]{\hsize}{\vskip7pt {\lantxh 胶胨罗斯氏菌\\可能导致败血症、脑膜炎、肺炎、菌血症等} \vskip7pt} & \parbox[c]{\hsize}{\vskip7pt\centerline{\raisebox{-1.5ex}{\includegraphics[width=5cm,height=0.55cm]{bar.pdf}}}\vskip7pt}  &
\hspace*{-4.83cm}\raisebox{-0.45ex}{\includegraphics{smile.pdf}}
 & \begin{minipage}{4.60cm}\begin{center}{{\lantxh 含量未超标{}} }\end{center} \end{minipage} \\
\hline
\parbox[c]{\hsize}{\vskip7pt {\lantxh 产气荚膜梭菌\\可导致气性坏疽，食物中毒} \vskip7pt} & \parbox[c]{\hsize}{\vskip7pt\centerline{\raisebox{-1.5ex}{\includegraphics[width=5cm,height=0.55cm]{bar.pdf}}}\vskip7pt}  &
\hspace*{-4.83cm}\raisebox{-0.45ex}{\includegraphics{smile.pdf}}
 & \begin{minipage}{4.60cm}\begin{center}{{\lantxh 含量未超标{}} }\end{center} \end{minipage} \\
\hline
\parbox[c]{\hsize}{\vskip7pt {\lantxh 病原性大肠埃希氏菌\\可能导致肠胃炎、尿路感染、新生儿脑膜炎、腹膜炎等} \vskip7pt} & \parbox[c]{\hsize}{\vskip7pt\centerline{\raisebox{-1.5ex}{\includegraphics[width=5cm,height=0.55cm]{bar.pdf}}}\vskip7pt}  &
\hspace*{-4.83cm}\raisebox{-0.45ex}{\includegraphics{smile.pdf}}
 & \begin{minipage}{4.60cm}\begin{center}{{\lantxh 含量未超标{}} }\end{center} \end{minipage} \\
\hline
\end{longtable}

%**************************
\vspace*{2mm}
\noindent
{\noindent\bf\xiaosihao 4、肠道菌营养功能}
\vspace*{1.5mm}

\begin{tctabularx}{tabularx={m{4.8cm}<{\centering}m{5.2cm}<{\centering}m{4.55cm}<{\centering}}}
&&
\\[-6pt]
\cellcolor{topcolor} \raisebox{2.614pt}{\color{black00}\bf 检测项目} &
\cellcolor{topcolor} \raisebox{2.614pt}{\color{black00}\bf 健康状态提示} &
\cellcolor{topcolor} \raisebox{2.614pt}{\color{black00}\bf 检测结果评价}
\end{tctabularx}

\vspace*{-4.25mm}
\fontsize{8.8pt}{11pt}\selectfont
\arrayrulecolor{gray2}
\begin{longtable}{m{4.8cm}m{5.2cm}<{\centering}m{0cm}@{}m{4.61cm}<{\centering}}
\hline
\parbox[c]{\hsize}{\vskip7pt {\lantxh 维生素B7\\维持皮肤及头发健康，增强免疫，与暗视力有关} \vskip7pt} & \parbox[c]{\hsize}{\vskip7pt\centerline{\raisebox{-1.5ex}{\includegraphics[width=5cm,height=0.55cm]{bar.pdf}}}\vskip7pt}  &
\hspace*{-3.17cm}\raisebox{-0.45ex}{\includegraphics{cry.pdf}}
 & \begin{minipage}{4.60cm}\begin{center}{{\lantxh 偏低{\\ \bahao 可能引发脱发和白发，不利于暗视力}} }\end{center} \end{minipage} \\
\hline
\multicolumn{3}{l}{}\\
\multicolumn{3}{l}{\includegraphics[scale=1]{greencircle.pdf} {\fontsize{9.3pt}{11pt}\selectfont\bf 绿色表示健康}   \includegraphics[scale=1]{orangecircle.pdf}  {\fontsize{9.3pt}{11pt}\selectfont\bf 橙色表示需要关注}  \includegraphics[scale=1]{redcircle.pdf} {\fontsize{9.3pt}{11pt}\selectfont\bf 红色表示有风险}}
\end{longtable}

\end{document}

